\chapter{METODOLOGIA}
\label{cap:metodologia}

\section{Protocolo e Relato da Revisão Sistemática}

Este trabalho adotou a metodologia de \textbf{Revisão Sistemática da Literatura}, fundamentada em princípios metodológicos consolidados do \textit{Cochrane Handbook for Systematic Reviews of Interventions} \cite{CochraneHandbook} e dos padrões do \textit{Institute of Medicine} (IOM). O relato dos resultados foi estruturado considerando as 27 recomendações do PRISMA 2020 (\textit{Preferred Reporting Items for Systematic Reviews and Meta-Analyses}) \cite{PRISMA2020}, buscando transparência, reprodutibilidade e completude na comunicação dos achados.

A escolha da abordagem PRISMA 2020 para estruturação do relato justificou-se por sua ampla aceitação na comunidade científica internacional como padrão de transparência e por orientar a documentação das etapas do processo de revisão. A replicabilidade depende, adicionalmente, dos parâmetros, registros e artefatos disponibilizados. Esta \textit{diretriz de relato} mostrou-se especialmente adequada para comunicar os resultados da Fase 1 do projeto, na qual o objetivo principal é mapear o estado da arte das técnicas computacionais aplicadas ao ensino de matemática.

\textbf{Importante:} Embora esta revisão não tenha sido registrada prospectivamente em bases como PROSPERO devido a restrições de cronograma acadêmico, o protocolo foi desenvolvido seguindo princípios de protocolos de revisões sistemáticas e do PRISMA-P 2015 (\textit{Preferred Reporting Items for Systematic Review and Meta-Analysis Protocols}) \cite{PRISMAP2015}, com definição \textit{a priori} de: (1) questão de pesquisa estruturada, (2) critérios de elegibilidade explícitos, (3) estratégia de busca reproduzível, (4) processo de seleção sistemático, e (5) método de síntese pré-definido. O checklist completo de aderência ao PRISMA 2020 encontra-se disponível no Apêndice~\ref{apendice:prisma-checklist}.

Não se afirma conformidade integral com o PRISMA-P; o protocolo é utilizado como referência para explicitar as decisões metodológicas adotadas.

\section{Estratégia de Busca}

\subsection{Bases de Dados e APIs}

A coleta de dados foi realizada mediante integração automatizada com as APIs de quatro bases científicas complementares:

\begin{enumerate}
    \item \textbf{Semantic Scholar}: Ampla cobertura em ciência da computação com métricas de influência científica;
    \item \textbf{OpenAlex}: Base de dados aberta e abrangente, sucessora do Microsoft Academic Graph;
    \item \textbf{Crossref}: Foco em metadados precisos de publicações e identificadores DOI;
    \item \textbf{CORE}: Agregador especializado em artigos de acesso aberto.
\end{enumerate}

A integração de múltiplas APIs de bases científicas \cite{CochraneHandbook} foi adotada para ampliar:

\begin{itemize}
    \item \textbf{Cobertura Complementar}: Cada base possui forças específicas --- \textit{Semantic Scholar} para métricas de impacto, \textit{OpenAlex} para amplitude de cobertura, \textit{Crossref} para precisão bibliográfica, e \textit{CORE} para acesso aberto;
    \item \textbf{Redução da dependência de uma única fonte}: amplia a cobertura consultada, sem eliminar vieses de seleção;
    \item \textbf{Rastreabilidade}: Automação via APIs e versionamento permitem repetir os passos e auditar parâmetros, embora as respostas externas possam variar;
    \item \textbf{Eficiência}: Coleta sistemática de grandes volumes de dados com consistência metodológica.
\end{itemize}

\subsection{Estratégia de Busca Bilíngue}

A estratégia de busca combinou três camadas de termos (base matemática, técnicas computacionais e domínio educacional), utilizando o operador booleano \texttt{AND} para buscar maior precisão e relevância temática.

\textbf{Camada 1 -- Base Matemática}:
\begin{itemize}
    \item Inglês: \textit{mathematics}, \textit{math} (2 termos)
    \item Português: \textit{matemática} (1 termo)
\end{itemize}

\textbf{Camada 2 -- Técnicas Computacionais} (12 termos para ambos os idiomas, com equivalentes em português):
\begin{itemize}
    \item \textit{adaptive} (adaptivo), \textit{personalized} (personalizado), \textit{tutoring} (tutor), \textit{analytics} (analítica), \textit{mining} (mineração), \textit{machine learning} (aprendizado de máquina), \textit{ai} (ia), \textit{assessment} (avaliação), \textit{student modeling} (modelagem do aluno), \textit{predictive} (preditivo), \textit{intelligent tutor} (tutor inteligente), \textit{artificial intelligence} (inteligência artificial)
\end{itemize}

\textbf{Camada 3 -- Domínio Educacional}:
\begin{itemize}
    \item Inglês: \textit{education}, \textit{learning} (2 termos)
    \item Português: \textit{educacao}, \textit{ensino} (2 termos)
\end{itemize}

A inclusão de termos de busca em português e inglês visou:

\begin{itemize}
    \item \textbf{Amplitude Geográfica}: Capturar pesquisas de diferentes regiões e contextos culturais;
    \item \textbf{Diversidade Cultural}: Incluir abordagens pedagógicas culturalmente específicas;
    \item \textbf{Ampliação de cobertura linguística}: aumentar a possibilidade de recuperar estudos relevantes em português e inglês, sem assegurar completude.
\end{itemize}

A estratégia resultou em \textbf{72 consultas únicas}:
\begin{itemize}
    \item 48 consultas em inglês (2 termos base $\times$ 12 técnicas $\times$ 2 educacionais)
    \item 24 consultas em português (1 termo base $\times$ 12 técnicas $\times$ 2 educacionais)
\end{itemize}

Cada consulta segue o formato: \texttt{"termo\_base" AND "termo\_tecnica" AND "termo\_educacional"}. 

\textbf{Exemplos}:
\begin{itemize}
    \item \texttt{"mathematics" AND "machine learning" AND "education"}
    \item \texttt{"math" AND "intelligent tutor" AND "learning"}
    \item \texttt{"matemática" AND "aprendizado de máquina" AND "educa\c{c}ao"}
\end{itemize}

Esta abordagem de \textbf{expansão em 3 camadas} busca oferecer:
\begin{enumerate}
    \item \textbf{Precisão temática}: Todas as queries combinam matemática + técnica computacional + domínio educacional
    \item \textbf{Cobertura de diferentes áreas}: os 12 termos técnicos contemplam diferentes áreas de IA/ML/LA
    \item \textbf{Inclusão bilinguística}: consultas em inglês e português podem ampliar a representação geográfica
    \item \textbf{Reprodutibilidade}: Estrutura documentada em \texttt{search\_terms.py} (módulo canônico)
\end{enumerate}

\section{Critérios de Seleção (PICOS)}

Os critérios de inclusão e exclusão foram definidos conforme o framework PICOS:

\begin{itemize}
    \item \textbf{P (Population)}: Estudantes de matemática em qualquer nível educacional;
    \item \textbf{I (Intervention)}: Aplicação de técnicas computacionais (ML, IA, LA, STI, NLP, etc.);
    \item \textbf{C (Comparison)}: Abordagens pedagógicas tradicionais ou alternativas (quando aplicável);
    \item \textbf{O (Outcomes)}: Desempenho acadêmico, diagnóstico de competências, personalização do ensino;
    \item \textbf{S (Study Design)}: Estudos empíricos, quasi-experimentais ou estudos de caso com evidências práticas.
\end{itemize}

\subsection{Justificativa do Recorte Temporal}

O período de 2015-2026 (com data de corte em 31 de agosto de 2026) foi escolhido por representar:

\begin{itemize}
    \item \textbf{Era da IA Educacional}: Década de maior evolução nas técnicas computacionais aplicadas à educação;
    \item \textbf{Maturidade do Machine Learning}: Consolidação de técnicas de ML em ambientes educacionais;
    \item \textbf{Explosão do Learning Analytics}: Desenvolvimento massivo de ferramentas de análise educacional;
    \item \textbf{Relevância Tecnológica}: Tecnologias ainda atuais e aplicáveis em contextos contemporâneos.
\end{itemize}

\subsection{Critérios de Inclusão}

\begin{enumerate}
    \item Artigos completos revisados por pares (\textit{peer-reviewed});
    \item Publicações entre 2015 e 2026 (com data de corte em 31 de agosto de 2026);
    \item Foco explícito em técnicas computacionais aplicadas ao ensino de matemática;
    \item Apresentação de dados empíricos, metodologias detalhadas ou evidências de desenvolvimento/avaliação de sistemas;
    \item Idiomas: inglês ou português.
\end{enumerate}

Esses itens registram os critérios planejados, mas não constituem uma confirmação bibliográfica individual de cada registro incluído. No snapshot atual há registros com veículo ausente, classificação institucional ou metadados ainda pendentes de confirmação; por isso, a seleção vigente não deve ser usada para afirmar que o status de texto completo e revisão por pares foi verificado para os 16 estudos.

\subsection{Critérios de Exclusão}

\begin{enumerate}
    \item Estudos com metodologia insuficiente ou incoerente;
    \item Trabalhos com foco indireto ou descontextualizado da matemática;
    \item Publicações predominantemente teóricas sem suporte empírico;
    \item Estudos com impacto não mensurável ou irrelevante;
    \item Documentos não validados cientificamente (\textit{preprints}, relatórios internos);
    \item Publicações com falhas conceituais ou contradições metodológicas;
    \item Idiomas diferentes de inglês ou português.
\end{enumerate}

\section{Processo de Seleção (PRISMA)}

O fluxo de seleção dos estudos foi organizado segundo as etapas PRISMA:

\subsection{Identificação}

Execução automatizada das consultas nas quatro APIs, resultando em \textbf{11.904 registros} no banco consolidado. Cada registro inclui metadados bibliográficos (título, autores, ano, \textit{venue}, DOI/URL) e, quando disponível, resumo (\textit{abstract}). A base \path{research/systematic_review.sqlite} foi adotada como fonte operacional local para os estados e contagens desta consolidação; ela não é um artefato versionado do repositório. O snapshot publicado é representado pelos exports versionados, pelo manifesto de reprodutibilidade e pelos arquivos bibliográficos e de auditoria associados. O snapshot foi persistido sem registros marcados pela flag \texttt{is\_duplicate}; por isso, não houve remoção adicional registrada no fluxo descrito neste capítulo. A auditoria de DOI é apresentada separadamente para não confundir a flag operacional com unicidade bibliográfica.

\subsection{Triagem}

Na triagem, os registros foram avaliados por regras de disponibilidade de título/resumo, período e idioma. Os mecanismos de identificação de possíveis duplicatas são descritos na Seção de Deduplicação e não devem ser confundidos com as exclusões de triagem. Foram aplicados filtros para remover:
\begin{itemize}
    \item Sem título ou resumo válido;
    \item Fora do período 2015--2026;
    \item Em idiomas não compatíveis.
\end{itemize}

No banco consolidado, os \textbf{11.904 registros} foram avaliados na triagem. \textbf{9.413 registros} foram excluídos nessa etapa, e \textbf{2.491 estudos} avançaram para a análise de elegibilidade.

\subsection{Elegibilidade}

Avaliação operacional dos registros mediante um sistema heurístico de pontuação (\textit{scoring}) de relevância temática, calculado sobre título e resumo quando disponíveis, baseado em:

\begin{enumerate}
    \item \textbf{Domínio matemático} (peso 0,3): ocorrência de termos de matemática e bônus para alguns tópicos específicos;
    \item \textbf{Relevância técnica} (peso 0,3): padrões de ML, IA, analítica da aprendizagem e algoritmos identificados no texto;
    \item \textbf{Sinais metodológicos} (peso 0,2): tipo de estudo e métodos de avaliação reconhecidos por padrões textuais;
    \item \textbf{Recência} (peso 0,1): faixa do ano de publicação;
    \item \textbf{Indicadores de metadados} (peso 0,1): comprimento do resumo, presença de DOI e indicador simplificado de veículo.
\end{enumerate}

A implementação soma os componentes ponderados, limita o resultado ao intervalo de 0 a 10 e, no snapshot atual, produziu valores observados entre 0 e 4,5. O limiar operacional foi definido como pontuação de relevância maior ou igual a 4,0. DOI, citações e acesso aberto não funcionam como medida de qualidade científica nessa implementação: apenas a presença de DOI e um indicador simplificado de veículo entram no componente de metadados. Por isso, o limiar não garante elegibilidade metodológica, qualidade dos estudos ou validade dos resultados, e os candidatos ainda requereram auditoria de pertinência.

Nesta etapa foram excluídos \textbf{2.475 registros} (99,4\% dos 2.491 registros avaliados). A pontuação foi um filtro operacional de aderência temática; essa contagem não constitui avaliação de qualidade metodológica.

Na reexecução e auditoria da rodada atual, foram identificados \textbf{23 candidatos} pelo limiar operacional. A revisão de pertinência identificou \textbf{7 falsos positivos}, que foram excluídos, permanecendo os \textbf{16 estudos} do conjunto vigente. Essa reconciliação explica a diferença em relação à execução histórica, que registrava 17 incluídos: a nova rodada utilizou o snapshot atualizado e a lógica de scoring corrigida, não uma simples substituição numérica. Os sete falsos positivos fazem parte dos 2.475 registros com estágio \texttt{eligibility} e status de exclusão; portanto, não são sete exclusões adicionais a serem somadas ao fluxo PRISMA.

\subsection{Inclusão}

Após a auditoria de pertinência dos candidatos acima do limiar, foram incluídos \textbf{16 estudos} para síntese qualitativa, representando uma taxa de inclusão de aproximadamente \textbf{0,13\%} em relação ao total identificado (16/11.904).

\section{Deduplicação}

É necessário distinguir a deduplicação prevista por helpers do módulo de processamento da lógica efetivamente usada na reexecução. O módulo legado oferece rotinas destrutivas por DOI, URL e similaridade de título; nesta última, TF-IDF (\textit{Term Frequency--Inverse Document Frequency}) com n-gramas de caracteres é comparado pela similaridade do cosseno, que varia de 0 a 1, com limiar operacional de $0{,}9$. Essas rotinas não devem ser descritas como se tivessem sido aplicadas ao snapshot atual.

\begin{enumerate}
    \item \textbf{Por DOI}: a rotina destrutiva normaliza o DOI (minúsculas e remoção do prefixo ``doi:'') e mantém a primeira ocorrência;
    \item \textbf{Por URL e similaridade de título}: as rotinas correspondentes podem remover repetições por URL e títulos com similaridade TF-IDF pelo cosseno $\geq 0{,}9$.
\end{enumerate}

Na execução canônica atual, o pipeline usa a identificação não destrutiva de possíveis duplicatas por DOI e URL: ela marca a flag, mas não remove linhas. Neste snapshot, os 11.904 registros possuem \texttt{is\_duplicate = 0} e o relatório registra zero remoções adicionais baseadas nessa flag. A auditoria do export encontrou 10.658 DOIs normalizados distintos, 1.222 registros sem DOI e 24 grupos de DOI repetido (48 linhas); esses casos requerem decisão de integridade própria e não foram convertidos automaticamente em exclusões. Assim, DOI repetido não equivale a duplicata removida, e nenhuma deduplicação por título deve ser inferida apenas a partir das contagens atuais.

\section{Infraestrutura Tecnológica}

\subsection{Pipeline Automatizado}

O pipeline de revisão sistemática foi implementado em \textbf{Python 3.11+}, utilizando:

\begin{itemize}
    \item \textbf{SQLite}: Banco de dados local para armazenamento estruturado (\path{systematic_review.sqlite});
    \item \textbf{Requests}: Cliente HTTP para comunicação com APIs;
    \item \textbf{Pandas}: Manipulação e análise de dados tabulares;
    \item \textbf{Scikit-learn}: Cálculo de similaridades TF-IDF para deduplicação.
\end{itemize}

O sistema implementa um \textit{cache} persistente na tabela \texttt{cache} do banco, indexado por API e consulta e com expiração configurada por fonte. Um acerto de cache ocorre quando a solicitação é respondida por um resultado previamente armazenado, sem nova chamada à API; o campo \texttt{hit\_count} é acumulado por entrada e os eventos de acerto/erro são registrados na auditoria. O cache reduz chamadas repetidas, mas não deduplica registros nem altera os estágios PRISMA. Sua taxa depende do histórico de execução e não deve ser tratada como uma constante do estudo.

\subsection{Reprodutibilidade}

O \textit{pipeline} foi projetado para ser reexecutável mediante comandos CLI (\textit{Command Line Interface}), com parâmetros, código e artefatos derivados versionados. Há, contudo, dois níveis de reprodução. A reconstituição do snapshot desta consolidação usa os exports versionados \path{research/exports/analysis/papers.csv}, \path{research/exports/analysis/papers.json}, o resumo PRISMA, a bibliografia dos estudos incluídos, o audit CSV e o manifesto \path{research/exports/reports/reproducibility_manifest.json}; nenhum desses artefatos exige distribuir o SQLite. A reexecução da coleta depende das respostas externas e do estado do cache, portanto pode produzir metadados e contagens diferentes.

O manifesto registra também os sete overrides aplicados após a auditoria de pertinência. Eles são necessários para reconstruir o conjunto vigente de 16 estudos, pois a regra de pontuação, isoladamente, ainda selecionaria os 23 candidatos operacionais. As justificativas substantivas individuais desses overrides permanecem uma pendência metodológica e não devem ser inferidas a partir do score.

Para reexecutar o pipeline em ambiente local, usando uma cópia do banco operacional quando disponível, podem ser usados:

\begin{verbatim}
# Verificar estatísticas de uma cópia local do banco operacional
python -m research.src.cli \
  --db /caminho/para/systematic_review.sqlite stats

# Regenerar exports a partir da cópia local
python -m research.src.cli \
  --db /caminho/para/systematic_review.sqlite export

# Atualizar o manifesto e seus hashes
python -m research.src.cli \
  --db /caminho/para/systematic_review.sqlite generate-manifest

# Executar nova coleta (pode divergir do snapshot publicado)
python -m research.src.cli run-pipeline --min-score 4.0

\end{verbatim}

Os parâmetros de configuração (limiares, APIs e termos de busca) estão documentados no repositório Git; credenciais e o banco operacional local permanecem fora do versionamento. A bibliografia completa do TCC é composta por \path{results/tcc/referencias.bib} e \path{results/tcc/referencias_pedagogicas.bib}. As referências de fundamentação teórica e metodológica não são contadas como estudos derivados do pipeline nem incluídas em \path{research/exports/references/included_papers.bib}.

\section{Avaliação da Qualidade Metodológica}

Foi definido um procedimento de avaliação crítica de qualidade metodológica utilizando o \textit{Mixed Methods Appraisal Tool} (MMAT) versão 2018 \cite{MMAT2018}. O MMAT foi selecionado por sua capacidade de avaliar diferentes desenhos de estudo (qualitativos, quantitativos e mistos) de forma padronizada, adequando-se à heterogeneidade metodológica identificada nesta revisão. Como a atualização do banco alterou o conjunto incluído, os julgamentos MMAT disponíveis para a versão anterior não são apresentados como avaliação final dos 16 estudos atuais; a reaplicação ao conjunto atualizado permanece pendente.

O processo de avaliação foi especificado nas seguintes etapas:

\begin{enumerate}
    \item \textbf{Classificação do desenho de estudo}: Os estudos incluídos foram categorizados conforme seu desenho metodológico (qualitativo, quantitativo randomizado, quantitativo não-randomizado, quantitativo descritivo ou métodos mistos).

    \item \textbf{Aplicação dos critérios MMAT}: Para cada estudo, deverão ser aplicados 5 critérios de qualidade específicos ao seu desenho metodológico, respondendo ``Sim'', ``Não'' ou ``Não é possível determinar'' para cada critério.

    \item \textbf{Síntese da qualidade}: Conforme recomendação dos autores do MMAT, a qualidade metodológica deverá ser reportada descritivamente (contagem de critérios atendidos) sem cálculo de \textit{score} percentual, evitando simplificação excessiva de aspectos complexos de qualidade.
\end{enumerate}

A avaliação planejada prevê registro por um único revisor, com justificativas detalhadas para cada julgamento, permitindo auditoria e revisão das decisões. Estudos com baixa qualidade metodológica não serão excluídos automaticamente, mas terão suas limitações consideradas na síntese narrativa e discussão dos achados.

O estado da avaliação MMAT do conjunto atualizado é registrado na Seção~\ref{sec:qualidade-estudos}, que explicita a pendência de reaplicação sem reutilizar a tabela histórica como resultado atual.

%--------- FIM METODOLOGIA ------------
